\section{Front-Loading Reasoning：推理能力是否应在预训练期打地基}

传统 pipeline 常把 pretraining 描述为世界知识学习，把 SFT 描述为格式模仿，把 RL 描述为真正的推理强化。Front-loading reasoning 质疑这种分工：如果 base model 从未在预训练中接触 reasoning trace，后续是否只能在较弱表示基础上修补？这里研究的不是“要不要 post-training”，而是固定预算下应该把多少推理数据提前。

\subsection{“Reasoning added too late”的弱地基问题}

课程用地基隐喻描述现有流程：pretraining 获得 facts，SFT 学会像 reasoning answer 的表面格式，RL/ RLVR 才鼓励真正探索。若 reasoning 只作为 post-hoc skill 加上去，base representation 可能缺少支持长链推断的早期结构。这个观点需要后续受控实验支持，不能只凭隐喻成立。

\lecturefigure{slide-014.jpg}{传统流程把 reasoning 放到后期，形成 weak foundation 的假设}{Stanford 官方录像，00:12:39--00:13:29}

\begin{knowledgebox}{背景概念：SFT 与 RLVR}
Supervised Fine-Tuning（SFT，监督微调）用给定输入--输出或 reasoning trace 做 teacher forcing，主要优化 likelihood；Reinforcement Learning with Verifiable Rewards（RLVR，可验证奖励强化学习）用可自动判定的答案、测试或规则给 rollout 奖励。二者都属于 post-training，但优化信号和数据要求不同。
\end{knowledgebox}

\subsection{实验轴：何时、多少、多广、多好}

Front-loading 的实验把同一类 reasoning data 分配到 pretraining、SFT 与后续阶段，并沿 diversity、quality、quantity 三个轴构造条件。No-reason base 在预训练中不看推理数据，reason base 则提前看到一部分。读图时先确认数据预算是否匹配，再看 SFT/RL 是否相同，否则“早加入”可能只是“总共加入更多”。

\lecturefigure{slide-015.jpg}{Front-loading reasoning 的阶段分配与三个数据轴}{Stanford 官方录像，00:13:29--00:14:54}

设总 reasoning token budget 为 $R$，其中 $R_{\mathrm{pre}}$ 分配到预训练，$R_{\mathrm{sft}}$ 分配到 SFT，则
\[
R_{\mathrm{pre}}+R_{\mathrm{sft}}=R.
\]
Front-loading 比较的是不同分配点，例如 $(0,R)$ 与 $(\alpha R,(1-\alpha)R)$。若总预算或后续训练轮数不同，必须另标为 compute-matched 或 data-matched 实验，不能当成同一问题。

\subsection{为什么 base 与 post-trained benchmark 不能混读}

只完成预训练的 base model 与完成 SFT/RL 的模型，通常使用不同 prompting、decoding 与 benchmark suite。图中把 General Purpose Reasoning、Math、Science、Code 等基准按阶段列开，目的是提醒读者：后面不同页的平均数未必来自同一组题，不能把 16\%、9.3\%、19\% 直接串成一条单调曲线。

\lecturefigure{slide-016.jpg}{Base model 与 post-trained model 使用的评测集合}{Stanford 官方录像，00:14:54--00:15:48}

\begin{warningbox}{读图时最容易犯的错误}
同名 benchmark 也可能因 few-shot、chain-of-thought、sampling temperature、pass@k、答案抽取器或 checkpoint 不同而不可直接比较。讲义保留课程的相对结论，但不会把跨阶段平均数解释成一个连续指标。
\end{warningbox}

